%! Author = lili
%! Date = 6/16/26


\section{Motivation and Background}\label{sec:motivation-and-background}
\begin{prob}
    \textbf{Incomplete Knowledge Graphs} \\
    Most existing \glspl{kg} suffer from limitations in knowledge
    sources.
    Biases during knowledge extraction lead to incomplete or erroneous information.
\end{prob}

\begin{bsp}
    A \gls{kg} about football might know ``Nunez plays for Liverpool'' but miss ``Nunez is a forward''
\end{bsp}

\defin{kgr}

\begin{application}
    \begin{itemize}
        \item Intelligent question-answering
        \begin{itemize}
            \item e.g.\ ``Who won the Champions League in 2020?''
        \end{itemize}
        \item Recommendation systems
        \begin{itemize}
            \item e.g.\ ``Users who liked X also liked Y''
        \end{itemize}
        \item Dialogue generation
        \begin{itemize}
            \item e.g.\ Chatbots with world knowledge
        \end{itemize}
    \end{itemize}
\end{application}

\begin{mytable}
    \begin{tabular}{p{2cm} p{5cm} p{5cm} }
       &  \textbf{Embedding-Based Methods} & \textbf{Logical Rule-Based Methods} \\
        Interpretability & Lack it & Strong \\
    TODO % TODO
    \end{tabular}\label{tab:table2}
\end{mytable}

\begin{keyin}
    Logical rule-based methods excel at inference by uncovering underlying logical rules,
    showcasing remarkable generalization ability and interpretability.
    They can be
    seamlessly integrated with \glspl{nn}.
\end{keyin}

\begin{figure}[H]
    % TODO 
    \caption{Taxonomy of Logical Rule-Based \gls{kgr} methods}\label{fig:figure}
\end{figure}

\begin{mytable}
    \begin{tabular}{ p{3cm} p{7cm} p{2cm}}
        \textbf{Category} & \textbf{Key Characteristics} & \textbf{Timeline} \\
        \gls{ilp}-Based & Inductive Logic Programming, Horn rules, rule
        learning from examples & 1990s \\
        Probabilistic Graphical
        Models + Rules & \gls{mln}, ProbLog, \gls{psl}, fuzzy reasoning & 2000s \\
        Embedding Techniques
        + Rules & \gls{kge}, rule em-
        bedding, mutual enhancement & 2010s \\
        \glspl{nn} + Rules & \gls{rnn}, \gls{gnn}, neural theorem prov-
        ing & 2018+
    \end{tabular}\label{tab:table3}
\end{mytable}
% TODO noch notizen hinzufügen?

\section{Background Knowledge}\label{sec:background-knowledge}

\defin{kg_formal}
\defin{fact_formal}

\begin{bsp} \defiref{fact_formal}
\[
    \langle Nunez, PlayForTeam, Liverpool \rangle
\]
\end{bsp}

% TODO triple types genauer recherchieren
% Triple Types: \gls{skg} (Static KG), \gls{tkg} (Temporal KG), MMKG (Multi-Modal KG)

\begin{figure}[H]
    % TODO  F 10
    \caption{An example of a knowledge graph showing entities (circles) and relations (edges)}\label{fig:figure2}
\end{figure}

\defin{kgdenoising}

% TODO F 11 generell

\begin{figure}[H]
    % TODO  F 11
    \caption{KGR task illustration}\label{fig:figure3}
\end{figure}

% TODO F 12 - 13

\section{\gls{ilp}-Based KGR Methods}\label{sec:ilp-based-kgr-methods}


\defin{ilp}

% TODO F 15 -17



\section{Probabilistic Graphical Models + Rules}\label{sec:probabilistic-graphical-models-+-rules}

The core idea of unifying probabilistic graphical models and logical rules is to combine the expressive power of \gls{fol}
with probabilistic reasoning to handle
uncertainty and incompleteness in \glspl{kg}.

\paragraph{Three Main Approaches:}
\defin{mln}
% TODO F 20
\defin{it_sld}
% TODO F 21
\defin{psl}
% TODO F 22

\begin{challenge}
    These methods rely on rule-based search mining requiring complex traversal
    computations and are therefore not well-suited for large-scale \glspl{kg}.
\end{challenge}



\section{Embedding Techniques + Rules}\label{sec:embedding-techniques-rules}

% TODO F 24 - 28


\section{Neural Networks + Rules}\label{sec:neural-networks-rules}

\subsection*{Why should we combine \glspl{nn} with rules?}
\glspl{nn} are fast, accurate, robust and good at generalization.
But they are limited due to being black-boxes and lacking explainability.
And our goal ist efficiency and explainable reasoning.

\begin{mytable}
    \begin{tabular}{ll}
        \textbf{Approach} & \textbf{Dominant Role} \\
        Rule-Enhanced \gls{nn} Modeling & \glspl{nn} (rules augment \gls{nn}) \\
        \gls{nn}-Enhanced Rule Learning & Rules (\glspl{nn} help learn rules) \\
    \end{tabular}
    \caption{The two subcategories of joining \glspl{nn} and logical rules.}
    \label{tab:nn_rules_roles}
\end{mytable}


\defin{lan}
\deffinS{law}

\begin{idee}
    \Glspl{lan} combine logical rule attention with neural attention.

    Final attention = weighted sum of $\alpha^{Logic}$ and $\alpha^{NN}$.

    This addresses query relation perception and neighborhood redundancy.
\end{idee}

\subsection*{AnyBURL}

TODO slide 32 % TODO

\subsection*{pLogicNet (\gls{gnn}-Based Approach)}
A Variational Expectation-
Maximization
Framework defines joint
distribution of all
possible triples using
\gls{mln} and \gls{fol}.

TODO slide 33 % TODO

\subsection*{VN Network: Virtual Neighbors for Emerging Entities}
TODO slide 34 % TODO

\subsection*{\gls{nn}-Enhanced Rule Learning}
\defin{ntp}
\Gls{ntp} ist based on the prolog programming language, uses a differentiable prover (DP) for differentiable proof
and recursive matching with \gls{nn} embeddings.

TODO slide 35 - 39
% TODO F 35 - 39

\section{Comparison and Analysis}\label{sec:comparison-and-analysis}

\begin{mytable}
    \begin{tabular}{llccc}
        \textbf{Category} & \textbf{Method} & \gls{mrr} & \textbf{\glssymbol{hitat}1} & \textbf{\glssymbol{hitat}10} \\

        \gls{ilp}-Based & AnyBURL & 0.310 & 0.233 & -- \\
        \gls{ilp}-Based & SAFRAN & 0.389 & 0.298 & 0.537 \\

        Embedding + Rules & RUGE & 0.768 & 0.703 & 0.815 \\
        Embedding + Rules & TARE & 0.781 & 0.721 & 0.839 \\

        \gls{nn} + Rules & VN & 0.463 & 0.345 & 0.526 \\
        \gls{nn} + Rules & GCR & 0.492 & -- & 0.490 \\
        \gls{nn} + Rules & GraIL & -- & -- & 0.641 \\
        \gls{nn} + Rules & CoMPILE & -- & -- & 0.847 \\
        \gls{nn} + Rules & ConGLR & 0.741 & -- & 0.830 \\
        \gls{nn} + Rules & BERTRL & 0.718 & -- & 0.867 \\
    \end{tabular}
    \caption{Performance Comparison on FB15K-237}
    \label{tab:kgr_performance}
\end{mytable}

\begin{mytable}
    \begin{tabular}{lcccc}
        \textbf{Category} & \textbf{Accuracy} & \textbf{Efficiency} & \textbf{Interpretability} & \textbf{Transferability} \\

        \gls{ilp}-Based &
        Low--Medium &
        Low--Medium &
            {\color{oceanmist}High} &
        Low \\

        Probabilistic + Rules &
        Low &
        Low &
        Medium &
        Low \\

        Embedding + Rules &
        Medium--{\color{oceanmist}High} &
            {\color{oceanmist}High} &
        Medium &
        Medium \\

        \gls{nn} + Rules &
            {\color{oceanmist}High} &
            {\color{oceanmist}High} &        Medium--{\color{oceanmist}} &            {\color{oceanmist}} \\
    \end{tabular}
    \caption{Comparative Analysis Summary.}
    \label{tab:kgr_comparison}
\end{mytable}

\begin{keyin}
    \begin{itemize}
        \item \gls{ilp}-based methods are most interpretable, but least efficient
        \item \gls{nn}+Rules results in the best accuracy and efficiency, but interpretability still challenging
        \item The tradeoff is Accuracy/Efficiency vs.
        Interpretability
    \end{itemize}
\end{keyin}

\section{Challenges and Future Directions}\label{sec:challenges-and-future-directions}
\begin{challenge}
    \textbf{Efficiency of Rule Learning} Path search complexity grows exponentially with KG size and effective
    sampling strategies or pre-trained models are needed.
\end{challenge}
\begin{challenge}
    \textbf{Rule Quality Evaluation} Traditional metrics (confidence, support) assume closed world (see \defirefshort{owa}). \gls{pca} is a compromise, not definitive solution and \gls{owa} prevalent in KGs makes missing data unusable
    as negative examples.
\end{challenge}
\begin{challenge}
    \textbf{Interpretability Balance} \gls{nn} integration improves speed but reduces interpretability.
    Embed
    reasoning steps in model architecture are needed.
\end{challenge}

\subsection*{Future Research Directions}
\begin{enumerate}
    \item \textbf{Better Benchmark Datasets}
    \begin{itemize}
        \item Current datasets (FB15K-237, NELL-995, WN18R) do not capture all logical patterns.
        \item Need datasets with pre-defined rules for proper evaluation.
    \end{itemize}

    \item \textbf{\gls{tkg} Reasoning}
    \begin{itemize}
        \item Uncover time-dependent or dynamic rules.
        \item Combine \gls{rnn}/LSTM with existing rule-learning methods.
    \end{itemize}

    \item \textbf{\gls{mmkg} Reasoning}
    \begin{itemize}
        \item Heterogeneous information (images, text, etc.).
        \item Use multi-modal learning for automatic rule extraction.
    \end{itemize}

    \item \textbf{Few-Shot and Zero-Shot Rule Discovery}
    \begin{itemize}
        \item Enhance representation with additional information.
        \item Leverage existing high-quality KGs as priors.
    \end{itemize}

    \item \textbf{Incorporating Logical Rules into LLMs (Large Language Models)}
    \begin{itemize}
        \item Address LLM hallucinations.
        \item Enhance pre-training, during-training, and post-training stages.
    \end{itemize}
\end{enumerate}

\subsection*{Key Takeaways}
\begin{enumerate}
    \item Knowledge graphs are often incomplete $\rightarrow$ reasoning is essential for filling gaps and correcting errors.

    \item Four main categories of logical rule-based KGR:
    \begin{itemize}
        \item \gls{ilp}
        \item Probabilistic + Rules
        \item Embedding + Rules
        \item \gls{nn} + Rules
    \end{itemize}

    \item Evolution:
    \begin{itemize}
        \item From interpretable but inefficient (\gls{ilp})
        \item To efficient but less interpretable (\gls{nn} + Rules)
    \end{itemize}

    \item Trade-off:
    \begin{itemize}
        \item Accuracy/Efficiency vs.
        Interpretability remains a key challenge.
    \end{itemize}

    \item Best performance is currently achieved by \gls{nn} + Rules hybrid methods.

    \item Future directions:
    \begin{itemize}
        \item Integration with LLMs
        \item \gls{tkg}/\gls{mmkg} reasoning
        \item Better evaluation benchmarks
    \end{itemize}
\end{enumerate}

